\documentclass[aspectratio=169, 11pt]{beamer}

% ── Configuración de idioma y fuentes ──
\usepackage[spanish, es-noshorthands]{babel}
\usepackage[utf8]{inputenc}
\usepackage[T1]{fontenc}
\usepackage{lmodern}
\usepackage{microtype}

% ── Configuración del tema Beamer ──
\usetheme{Madrid}
\usecolortheme{default}

% ── Colores corporativos (SHCP - UPIT / VINO) ──
\definecolor{vino}{HTML}{6E152E}
\definecolor{vinoDark}{HTML}{4B1A28}
\definecolor{vinoLite}{HTML}{b84a6b}

% Aplicar colores al tema de la presentación
\setbeamercolor{palette primary}{bg=vino,fg=white}
\setbeamercolor{palette secondary}{bg=vinoDark,fg=white}
\setbeamercolor{palette tertiary}{bg=vinoLite,fg=white}
\setbeamercolor{structure}{fg=vino}
\setbeamercolor{title}{bg=vino,fg=white}
\setbeamercolor{frametitle}{bg=vino,fg=white}
\setbeamercolor{item projected}{bg=vino,fg=white}

% Eliminar iconos de navegación abajo
\setbeamertemplate{navigation symbols}{}

% ── Datos de la Portada ──
\title[Conciliador v6.1]{Conciliador de Cotizaciones}
\subtitle{Solución automatizada de extracción y análisis de presupuestos}
\author[SHCP - UPIT]{SHCP - UPIT}
\date{\today}

\begin{document}

% ── Diapositiva 1: Portada ──
\begin{frame}
    \titlepage
\end{frame}

% ── Diapositiva 2: El Problema ──
\begin{frame}{El Problema Actual}
    \begin{itemize}[<+->]
        \item \textbf{Procesos Manuales Lentos:} La transcripción manual de datos económicos desde archivos PDF a hojas de Excel consume mucho tiempo.
        \item \textbf{Riesgo de Errores:} El factor humano en el tipeo de cifras largas (presupuestos de obra, tecnología, etc.) genera discrepancias millonarias.
        \item \textbf{Falta de Estandarización:} Cada proveedor utiliza formatos distintos para presentar subtotales, IVA y tipos de moneda.
        \item \textbf{Cálculos Cambiarios Tediosos:} Las conversiones entre USD, EUR y MXN requieren revisar manualmente el tipo de cambio oficial al día.
    \end{itemize}
\end{frame}

% ── Diapositiva 3: Nuestra Solución ──
\begin{frame}{Nuestra Solución: Conciliador v6.1}
    Desarrollamos una \textbf{aplicación de escritorio nativa} (PyQt6) robusta y autosuficiente.
    
    \vspace{0.5cm}
    \begin{columns}
        \column{0.5\textwidth}
        \textbf{¿Qué logra?}
        \begin{itemize}
            \item Lee archivos PDF automáticamente.
            \item Extrae cantidades, rubros e IVA de forma inteligente.
            \item Agrupa varias cotizaciones en un único flujo de trabajo.
            \item Genera reportes en Excel listos para la conciliación manual final.
        \end{itemize}
        
        \column{0.5\textwidth}
        \begin{center}
            \colorbox{vino!10}{
                \parbox{0.9\textwidth}{
                \centering \textbf{\color{vino}Impacto Directo:}\\
                Reducción del 80\% del tiempo de transcripción y 0\% de errores tipográficos en la extracción.
                }
            }
        \end{center}
    \end{columns}
\end{frame}

% ── Diapositiva 4: Características Principales ──
\begin{frame}{Características Técnicas Principales}
    \begin{enumerate}[<+->]
        \item \textbf{5 Estrategias de Análisis PDF:} Búsqueda de tablas estructuradas, barrido por texto, extracción geométrica de palabras clave y Motor OCR (Reconocimiento Óptico) opcional.
        \item \textbf{API Banxico Integrada:} Consulta automática del tipo de cambio oficial al día (Serie SF43718, etc.) vía REST API sin salir de la app.
        \item \textbf{Auditoría y Trazabilidad:} Calcula un código MD5 único (Hash) que fusiona y sella digitalmente los PDFs para probar su integridad y autenticidad.
        \item \textbf{Exportación Inteligente:} Plantillas Excel pre-configuradas con diseño corporativo y fórmulas integradas listas para imprimir (Tamaño Carta, Vertical).
    \end{enumerate}
\end{frame}

% ── Diapositiva 5: Demostración de Interfaz ──
\begin{frame}{Flujo de Trabajo (Interfaz)}
    \begin{center}
        \colorbox{vino}{\color{white} 1. Cargar PDFs} $\rightarrow$
        \colorbox{vino}{\color{white} 2. Seleccionar Páginas} $\rightarrow$
        \colorbox{vino}{\color{white} 3. Extraer Montos} $\rightarrow$
        \colorbox{vino}{\color{white} 4. Validar} $\rightarrow$
        \colorbox{vino}{\color{white} 5. Exportar}
    \end{center}
    \vspace{0.4cm}
    \begin{itemize}
        \item \textbf{Visor Integrado:} Permite hojear los PDFs lado a lado con los datos extraídos.
        \item \textbf{Recálculo en Vivo:} La tabla recalcula el IVA y el Subtotal en tiempo real si el usuario decide modificar alguna cantidad manualmente.
        \item \textbf{Gestión de Memoria Segura:} Agrega y quita secciones de proveedores sin perder datos que ya fueron extraídos.
    \end{itemize}
\end{frame}

% ── Diapositiva 6: Cierre ──
\begin{frame}
    \begin{center}
        \Huge \textcolor{vino}{\textbf{¡Gracias!}}
        
        \vspace{1cm}
        \Large \textbf{¿Tienen alguna pregunta?}
        
        \vspace{1cm}
        \normalsize Secretaría de Hacienda y Crédito Público (SHCP)\\
        Unidad de Política e Ingresos Tributarios (UPIT)
    \end{center}
\end{frame}

\end{document}
